\documentclass[10pt,a4paper]{amsart}
\usepackage[margin=19mm]{geometry}
\usepackage{amsmath,amssymb,mathtools}
\usepackage[scr=rsfs,cal=euler]{mathalfa}
\usepackage{xfrac}
\usepackage[unicode,hidelinks,bookmarks=false]{hyperref}
\newcommand{\ie}{i.e.\ }
\newcommand{\cf}{cf.\ }
\newcommand{\resp}{respectively\ }
\newcommand{\Subsection}{\subsection}
\providecommand{\nobreakdash}{\nobreak\hbox{-}}
\setlength{\emergencystretch}{2em}
\multlinegap0pt
\usepackage{bm}
\usepackage{bbm}
\usepackage{dsfont}
\usepackage{xspace}
\usepackage{cancel}
\usepackage{mathtools}
\usepackage[shortlabels]{enumitem}
\usepackage{amsthm}
\makeatletter
\@ifundefined{definition}{\newtheorem{definition}{Definition}}{}
\@ifundefined{lemma}{\newtheorem{lemma}{Lemma}}{}
\@ifundefined{theorem}{\newtheorem{theorem}{Theorem}}{}
\makeatother
\renewcommand{\hat}{\widehat}
\renewcommand{\tilde}{\widetilde}
\title[A Generating Polynomial Based Two-Stage Optimization Method ]{A Generating Polynomial Based Two-Stage Optimization Method for Tensor Rank Decomposition}
\author{Zequn Zheng, Hongchao Zhang, Guangming Zhou}
\date{Source: arXiv:2504.00313v3 (CC BY 4.0)}
\begin{document}
\maketitle

\noindent\emph{Source and licence.} A Generating Polynomial Based Two-Stage Optimization Method for Tensor Rank Decomposition. Version arXiv:2504.00313v3 (CC BY 4.0), distributed under the Creative Commons Attribution 4.0 International licence, as recorded in the arXiv licence metadata for this exact version and shown on the abstract page https://arxiv.org/abs/2504.00313v3. Full rights notice: licenses/arxiv-2504.00313-CC-BY-4.0.txt.\par\smallskip

\noindent\textit{Editorial scope.} Counted unit: the Theorem whose source label is \texttt{thm:SequDecom} in the article (source0.tex lines 675--687, source label \texttt{thm:SequDecom}), delivered together with the article's own complete original proof of it (source0.tex lines 688--744, one \texttt{proof} environment of 57 lines). No mathematics is written, repaired or re-derived: statement and proof are the article's own text line for line. Required context retained uncounted: lemma:rton1 (lines 484--491); eq:pre-processing (lines 555--557); T-decomp (lines 565--567); def-gen-eigen (lines 654--662). Every definition, display and label of those lines is retained unchanged. Omitted from this package and not redistributed: 1--483, 492--554, 558--564, 568--653, 663--674, 745--1877. Nothing inside a retained excerpt is omitted or altered: every retained body is delivered exactly as the article's lines are written, so no omission receipt of any kind accompanies this package. Citations: the retained excerpts carry no citation command at all, so no bibliography entry of the article is transcribed and no reference is redistributed. Reference adaptation: none was needed and none was made; every reference inside the delivered unit resolves inside it, so no unresolved reference or citation is produced. Printed slips kept literal: no printed word, symbol, hypothesis or number of the counted statement or of its proof was altered. Layout and class: the article's own class is \texttt{siamltex} with packages including graphicx, color, xcolor, bm, amssymb, amsfonts, amsmath, mathtools, enumerate, url; the wrapper is the lane's pinned \texttt{amsart} editorial wrapper at 10pt on A4, which restates the theorem-like environments the retained text needs and adds the article's own macro definitions that the retained text needs (\texttt{\textbackslash hat}, \texttt{\textbackslash tilde}), with extra wrapper packages (bm, bbm, dsfont, xspace, cancel, mathtools, enumitem). The delivered numbering is the unit's own. Assets: no retained span contains a figure environment, a graphics-inclusion command, a PDF-page inclusion, a table or a TikZ picture; no image file is copied into this package, the package declares no asset dependency and no image file is redistributed with it. Licence: version arXiv:2504.00313v3 of this article is distributed under the Creative Commons Attribution 4.0 International licence, as recorded in the arXiv licence metadata for this exact version and shown on the abstract page https://arxiv.org/abs/2504.00313v3. A full rights notice is delivered at licenses/arxiv-2504.00313-CC-BY-4.0.txt. Characters: every retained excerpt is pure ASCII, so the delivered engine, XeLaTeX with its default Latin Modern text font, reports no missing character.
\begin{lemma} \label{lemma:rton1}
Let $\mathcal{F}\in \mathbb{C}^{n_1 \times n_2 \times  n_3}$ be an order-3 tensor with
 rank $n_2 < r \le n_1$. 
 Suppose $\mathcal{F}_{1:r,:,:}=\hat{U}^{(1)} \circ \hat{U}^{(2)} \circ \hat{U}^{(3)}$
 and let  $A^1= \hat{U}^{(2)} \odot \hat{U}^{(3)}$ and $B^1=\mathrm{Flatten}(\mathcal{F},1)^\top$.
Then, in the generic case, the linear system  $A^1X=B^1$ has a least squares solution
 $\tilde{U}^{(1)}$ and  $\mathcal{F}=\tilde{U}^{(1)} \circ \hat{U}^{(2)}  \circ \hat{U}^{(3)}$.
\end{lemma}

\begin{equation}\label{eq:pre-processing}
    P \mathcal{F}_{1:r,:,1}=(I_r)_{:,1:n_2}.
\end{equation}

\begin{equation}\label{T-decomp}
 \mathcal{T} = \hat{U}^{(1)} \circ U^{(2)} \circ U^{(3)}, 
\end{equation}

\begin{definition} \label{def-gen-eigen}
    For a set of matrices $A_1,\cdots,A_d \in \mathbb{C}^{m \times n}$ with $m \geq n$, a full rank matrix $S \in \mathbb{C}^{m \times m}$ is called the generalized left common eigenmatrix of $A_1,\cdots,A_d$, if it satisfies
    \begin{equation} \label{eq:gceig}
            S A_k  = D_k S_{:,1:n} ~~\mathrm{for~}1\leq k \leq d, 
    \end{equation}
    where $D_k \in \mathbb{C}^{m \times m}$ is a diagonal matrix. 
    Then, for all $i =1, \ldots,m$, $s^i \coloneqq S_{i,:}$ is called a generalized left common eigenvector of 
$A_1,\cdots,A_d$, and $\lambda_{i,k} \coloneqq (D_k)_{i,i}$s is called the generalized left common eigenvalue of $A_k$ associated with $s^i$.
\end{definition}

\begin{theorem} \label{thm:SequDecom}
Let $\mathcal{F}\in \mathbb{C}^{n_1 \times n_2 \times  n_3}$ be an order-3 tensor with
 rank $n_2 < r \le n_1$. 
 Suppose $\mathcal{T}$ is the reduced tensor of $\mathcal{F}$ given in \eqref{eq:pre-processing}
 with  ${T}_{k}={\mathcal{T}}_{:,:,k}$, $k=1, \ldots, n_3$.
In the generic case, we have
    \begin{enumerate}[(i)]
        \item for each nonsingular generalized left common eigenmatrix $S$ of $T_2,\cdots,T_{n_3}$,
         $\mathcal{F}$ has a tensor decomposition given in \eqref{eq:findTDbygceig};
        \item  for each tensor decomposition, there is a nonsingular generalized left 
        common eigenmatrix $S$ of $T_2,\cdots,T_{n_3}$.
    \end{enumerate}
\end{theorem}

\begin{proof}
We first prove $(i)$. Suppose $S$ is a nonsingular generalized left common eigenmatrix of 
$T_2,\cdots,T_{n_3}$, that is 
\begin{equation}\label{S-Tk-1}
    S T_k = D_k S_{:,1:n_2}~~\mathrm{for~} 2\leq k \leq n_3. 
\end{equation}
Let $\lambda_{i,k}=(D_k)_{i,i}$ for all $1\leq i \leq r$ and $2\leq k \leq n_3$, and
let $\hat{\mathcal{T}} = \hat{U}^{(1)} \circ U^{(2)} \circ U^{(3)}$, where
\begin{align}\label{eq:decombyS}
    \hat{U}^{(1)} = S^{-1},~U^{(2)} = (S_{:,1:n_2})^\top \mbox{ and } U^{(3)} = \begin{pmatrix}
        1 & 1 & \dots & 1 \\
        \lambda_{1,2} & \lambda_{2,2} & \dots & \lambda_{r,2} \\
        \vdots  &\vdots  &\vdots  &\vdots  \\
         \lambda_{1,n_3} & \lambda_{2,n_3} & \dots & \lambda_{r,n_3} \\
    \end{pmatrix}.
\end{align}
Then, by the construction of $\hat{\mathcal{T}}$ and \eqref{S-Tk-1}, 
\begin{align*}
\hat{\mathcal{T}}_{:,:,k}=\hat{U}^{(1)}\mathrm{diag}(U^{(3)}_{k,:})(U^{(2)})^\top
=S^{-1}D_{k}S_{:,1:n_2} = S^{-1} (S T_k) =T_k.
\end{align*}
Hence, we have $\hat{\mathcal{T}}=\mathcal{T}$. 
Then, it follows from \eqref{eq:pre-processing} that
\begin{align*}
 \mathcal{F}_{1:r,:,:}=P^{-1}\times_1 \mathcal{T}= P^{-1}\times_1 \hat{\mathcal{T}} =
 P^{-1} \hat{U}^{(1)} \circ U^{(2)} \circ U^{(3)}.
\end{align*}
This gives a tensor decomposition for $\mathcal{F}_{1:r,:,:}$.
Then, by Lemma \ref{lemma:rton1}, the linear least squares system $AX=B$
has a solution, denoted as ${U}^{(1)}$, where $A=U^{(2)} \odot U^{(3)}$, 
$B=\mathrm{Flatten}(\mathcal{F},1)^\top$, and we have a tensor decomposition of 
$\mathcal{F}$ as 
\begin{align} \label{eq:findTDbygceig}
    \mathcal{F}={U}^{(1)} \circ U^{(2)} \circ U^{(3)}.
\end{align}

We now prove $(ii)$. 
This essentially follows from the previous discussion on the motivations of the Definition~\ref{def-gen-eigen}.
Since $\mathcal{T}$ is the reduced tensor of $\mathcal{F}$ given in \eqref{eq:pre-processing}, we have \eqref{T-decomp} holds. That is 
$\mathcal{T} = \hat{U}^{(1)} \circ U^{(2)} \circ U^{(3)}$, 
where $\hat{U}^{(1)} = PU^{(1)}_{1:r,:}$, $P$ is given in \eqref{eq:pre-processing}
and $ U^{(i)} $, $i=1,2,3$, are matrices such that 
$\mathcal{F} = U^{(1)} \circ U^{(2)} \circ U^{(3)}$.
Let $S = (\hat{U}^{(1)})^{-1}$.
Then, for all $1 \leq k \leq n_3$, we have
\begin{equation}\label{TkTkTk}
T_k= PU^{(1)}_{1:r,:} \text{diag}(U^{(3)}_{k,:})({U}^{(2)})^\top=S^{-1} \text{diag}(U^{(3)}_{k,:})({U}^{(2)})^\top.
\end{equation}
It then follows from  $T_1=(I_r)_{:,1:n_2} $,  
$\mathrm{diag}(U^{(3)}_{1,:}) = I_r$ and \eqref{TkTkTk} that $S_{:,1:n_2} =  (U^{(2)})^\top$.
Hence, by  \eqref{TkTkTk}, for $2 \leq k \leq n_3$ we have
\begin{align}
    S T_k = \text{diag}(U^{(3)}_{k,:})({U}^{(2)})^\top= \text{diag}(U^{(3)}_{k,:}) S_{:,1:n_2}.
\end{align}
Therefore, $S$ is a nonsingular generalized left common eigenmatrix of 
$T_2,\cdots,T_{n_3}$.
\end{proof}

\end{document}
